\documentclass{article}
\usepackage[UTF8]{ctex}
\usepackage{amsmath}
\usepackage{amssymb}
\usepackage{graphicx}
\usepackage{geometry}
\geometry{a4paper, margin=2cm}

\newcommand{\problemname}[1]{\begin{center}\Large\bfseries #1\end{center}}
\usepackage{enumitem}

\begin{document}
\problemname{蒜头君当裁判}

蒜头君在某个比赛当中当裁判。比赛规则比较奇怪，要求选手写出一个两位整数，如果选手写出一个个位和十位数字一样的整数，如 $55$，那么就判定他胜利了，此时输出 \texttt{win}，否则输出 \texttt{lose}。

现在蒜头君拿到了一张写着一个两位整数的纸条，我们跟他一起来当裁判吧。

\subsection*{输入格式}
一个整数 $n\text{ }(10\len\le99)$，表示蒜头君拿到的纸条上写的整数。

\subsection*{输出格式}
按照规则输出 \texttt{win} 或 \texttt{lose}。

\subsection*{样例输入}
\begin{verbatim}
77
\end{verbatim}

\subsection*{样例输出}
\begin{verbatim}
win
\end{verbatim}

\subsection*{样例 2 输入}
\begin{verbatim}
10
\end{verbatim}

\subsection*{样例输出}
\begin{verbatim}
lose
\end{verbatim}


\end{document}
